\subsection{Comparing and Harmonizing Regional AI Rules}
\label{sec:8.7.4}
Four governments have made four different bets on how to regulate AI within their own borders: the European Union a comprehensive risk-tiered statute, the United States a sectoral patchwork built on agencies that predate AI, the United Kingdom principles routed through existing regulators, and China a licensing regime built to police a communications channel rather than a product (section~\ref{sec:8.7.8} takes each in turn). Those four already diverge sharply. The regions outside those four diverge again, in different directions, and what actually crosses a border between any of them turns out to be narrower than a shared statute.

Canada regulates its own federal government's use of AI through a binding rule, the 2019 Directive on Automated Decision-Making, which requires an algorithmic impact assessment, scaled to the decision's stakes, before a federal system can make or assist a decision that affects a person \autocite{tbs2019automated}. A broader statute meant to extend comparable obligations to private companies, the Artificial Intelligence and Data Act, spent three years moving through Parliament and then died with it: when Parliament was prorogued in January 2025, the underlying bill, C-27, died on the order paper, and had not been reintroduced as of mid-2026. A framework that depends on one government's survival is a fragile thing for a trading partner to harmonize with.

Singapore took the opposite structural bet: no statute, a voluntary Model AI Governance Framework that its privacy regulator and infocomm authority first published in 2019 and revised in 2020, giving companies detailed but non-binding guidance on internal AI governance \autocite{pdpc2020model}. Since 2022, a companion toolkit called AI Verify has let a company test its own system against the framework's principles and produce a report a regulator, a customer, or the public can inspect. It is the UK's wager taken further: not just guidance routed through an existing regulator, but guidance backed by no regulator's legal mandate at all, only the market pressure of a public test result.

Taiwan took a third path, and only recently. Its Legislative Yuan passed an AI Basic Act in December 2025, promulgated the following month as the country's first dedicated AI statute: twenty articles built around seven principles, among them transparency and explainability, fairness and non-discrimination, and accountability, administered by the National Science and Technology Council \autocite{moda2025basic}. The principles echo the UK's five almost point for point. What differs is the instrument: Taiwan wrote them into binding statute rather than administrative guidance, landing somewhere between the EU's comprehensive-statute model and the UK's lighter one.

One place harmonization is actually happening, rather than being proposed, is narrower than any of this: cross-border data flows. In March 2022, the US Department of Commerce and Singapore's trade and infocomm authorities agreed to align their AI governance frameworks and expand Singapore's participation in APEC's Cross-Border Privacy Rules System, a certification that lets a company demonstrate it meets one common privacy standard and have every participating economy recognize that single certification rather than auditing the company separately in each \autocite{commerce2022singapore}. That is what cross-border harmonization actually looks like in practice: not a shared global framework, but a small set of governments willing to certify to the same standard and recognize each other's certifications.

None of this converges toward a single global standard, and nothing here suggests it will. What actually crosses a border is a certification a handful of regulators agree to recognize, a technical standard a company adopts because a customer requires it, or a treaty obligation written at the level of human rights rather than a specific rule — section~\ref{sec:8.7.2} covers the first two mechanisms, and section~\ref{sec:8.7.8} covers the third. Harmonizing regional AI rules is less a matter of writing one shared framework than of finding, case by case, a mechanism a government working from a different premise still has reason to follow.
